\documentclass[11pt, a4paper]{article}
\usepackage[a4paper, top=2.5cm, bottom=2.5cm, left=2cm, right=2cm]{geometry}
\usepackage{amsmath,amssymb,amsthm,microtype}
\usepackage{hyperref}

\theoremstyle{definition}
\newtheorem{problem}{Problem}
\theoremstyle{plain}
\newtheorem{theorem}{Theorem}
\newtheorem{lemma}{Lemma}

\begin{document}

\title{\vspace{-1.5cm}\textbf{Regional Quantitative Problem-Solving Circle} \\ \Large Problem Set 6: Stochastic Modeling, Markov Chains \& Synthesis (Weeks 11--12)}
\author{\textbf{Lead Architect:} Mohamed Jad Srifi}
\date{\textbf{Term:} Fall 2025 -- Spring 2026}
\maketitle

\noindent\textbf{Instructions:} Formalize all random processes by defining their state spaces, transition probabilities, and boundary conditions. Use First-Step Analysis to construct and resolve difference equations or linear recurrences.

\vspace{0.5cm}

\begin{problem}[Gambler's Ruin / First-Step Difference Equations --- Putnam Caliber]
A gambler starts with an initial fortune of $\$k$ and bets $\$1$ repeatedly on independent coin flips. The probability of winning a given flip is $p$, and the probability of losing is $q = 1 - p$. The game strictly terminates when the gambler's fortune reaches $\$0$ (ruin) or $\$N$ (bankrupting the house).
Define $R_k$ as the probability of ultimate ruin given a starting state of $k$.
Derive the boundary value recurrence relation for $R_k$. By analyzing the telescoping differences $(R_{k+1} - R_k)$, derive the exact closed-form expressions for $R_k$ in the symmetric case ($p = 1/2$) and the asymmetric case ($p \ne 1/2$).
\end{problem}

\vspace{0.5cm}

\begin{problem}[Expected Hitting Times on Symmetric Cyclic Graphs --- MIT 6.042 Standard]
A particle performs a simple random walk on a finite cycle consisting of $N$ nodes consecutively labeled $0, 1, \dots, N-1$. At each discrete time step, the particle transitions to its clockwise neighbor with probability $1/2$ or to its counter-clockwise neighbor with probability $1/2$. Node $0$ is designated as an absorbing boundary state.
Let $E_k$ be the expected number of time steps required to reach node $0$, given that the particle originates at node $k$.
Derive the system of linear equations for $E_k$, establish the second-order finite difference invariant $\Delta^2 E_k$, and rigorously prove that $E_k = k(N - k)$ for all $0 \le k \le N-1$.
\end{problem}

\vspace{0.5cm}

\begin{problem}[Synthesis: Coupling \& Stationary Distribution Convergence]
Two independent particles, $A$ and $B$, execute lazy random walks on the complete graph $K_M$ (where $M \ge 3$). At each discrete time step, each particle independently chooses to remain at its current vertex with probability $1/2$, or transitions to a uniformly chosen distinct vertex with probability $\frac{1}{2(M-1)}$.
Formulate this dual-particle system as a single absorbing Markov chain based on spatial intersection. Calculate the exact expected number of steps required for the two particles to simultaneously occupy the exact same vertex, assuming they originate on distinct vertices.
\end{problem}

\end{document}
